% pdflatex-tate-notes.tex -- T3: tate 脚注/後注/傍注フィクスチャ
% 検証内容（uplatex+tate を ground truth とした実測比較で確認済みの項目：
% 脚注=最左端カラム上揃え／後注=本文直後のカラム／傍注(marginpar)=本文カラムの
% 下方、いずれも縦組み正立，数ミリの家具由来オフセットを除き uplatex と一致）:
%   NOTE-JAREG   : \footnotesize で \jlreq@jareg(\zw実寸)が 8pt に追従する
%                  （横組 (a) バグの縦組挿入物内回帰ガード．旧バグでは101pt）
%   NOTE-FNCOUNT : footnote/endnote/sidenote が共有カウンタを正しい順で
%                  進める（本フィクスチャの系列で 7；uplatex GT でも同一）
%   NOTE-THEFN   : 縦組の \thefootnote（4.2.2注3, （＊）形式）が展開可能で
%                  空でない（\tatechuyoko* は pdfTeX では既知のパススルー）
%   NOTE-ENDFLUSH: \theendnotes 出力で後注リストが空に戻る
%   コンパイル成立自体が \sidenote/\sidenotemark/\sidenotetext/
%   \marginpar/複数ページ脚注(twoside 遅延経路は別途 stress 確認)のゲート
\documentclass[pdflatex,tate,paper=a5,sidenote_length=15zw]{jlreq}
\pagestyle{empty}
\begin{document}
\makeatletter
% --- NOTE-JAREG: 挿入物箱内のフォントサイズ追従
\setbox\@tempboxa=\vbox{\reset@font\footnotesize
  \xdef\jlreq@testjareg{\the\jlreq@jareg}X}%
\ifdim\jlreq@testjareg=\dimexpr.8\jlreq@fontsize\relax
  \typeout{TATE-PASS: NOTE-JAREG (\jlreq@testjareg)}%
\else
  \typeout{TATE-FAIL: NOTE-JAREG (\jlreq@testjareg, expected \the\dimexpr.8\jlreq@fontsize\relax)}%
\fi
\makeatother
% --- 脚注（短いもの + 分離型）
本文に脚注を付ける\footnote{縦組みの脚注テキスト。}さらに続く本文。
分離型も\footnotemark 使う\footnotetext{分離した脚注。}。
% --- 長い脚注（禁則ストレス：「」や、。が行をまたぐ位置に来る長さ）
禁則の確認を兼ねた長い脚注\footnote{これは禁則処理の確認のための長めの脚注です。
「開き括弧や句読点、そして閉じ括弧」が行頭や行末に来た場合の処理を確認するための
テキストを含んでいます。さらに長くするために文を足します。縦組みの脚注は
欄外側の左端に、上揃えで配置されるべきものです。}。
% --- 後注 ×2（カウンタ共有の確認用）
後注を一つ使う\endnote{縦組みの後注テキスト一。}。
% --- 傍注（sidenote_type=number 既定）と単体 marginpar
傍注の呼び出し\sidenote{縦組みの傍注テキスト。}と、
分離型\sidenotemark も使い\sidenotetext{分離した傍注。}、
単体の欄外メモ\marginpar{欄外メモ「括弧」つき。}も併記する。
後注をもう一つ\endnote{縦組みの後注テキスト二。}。
% --- NOTE-FNCOUNT / NOTE-THEFN / NOTE-ENDFLUSH
\makeatletter
% カウンタ共有経路（uplatex GT と同じ歩み）: \footnote→1, \footnotemark→2,
% 長い\footnote→3, \endnote一→4, \sidenote→5, \sidenotemark→6, \endnote二→7
\ifnum\c@footnote=7\relax
  \typeout{TATE-PASS: NOTE-FNCOUNT (\the\c@footnote)}%
\else
  \typeout{TATE-FAIL: NOTE-FNCOUNT (\the\c@footnote, expected 5)}%
\fi
\protected@edef\jlreq@testfn{\thefootnote}%
\ifx\jlreq@testfn\@empty
  \typeout{TATE-FAIL: NOTE-THEFN (empty)}%
\else
  \typeout{TATE-PASS: NOTE-THEFN (\jlreq@testfn)}%
\fi
\ifx\jlreq@endnote@texts\@empty
  \typeout{TATE-FAIL: NOTE-ENDPRE (endnote list empty before \string\theendnotes)}%
\else
  \typeout{TATE-PASS: NOTE-ENDPRE (endnotes pending)}%
\fi
\makeatother
\theendnotes
\makeatletter
\ifx\jlreq@endnote@texts\@empty
  \typeout{TATE-PASS: NOTE-ENDFLUSH (list flushed)}%
\else
  \typeout{TATE-FAIL: NOTE-ENDFLUSH (list not flushed)}%
\fi
\makeatother
\newpage
二ページ目の本文\footnote{二ページ目の脚注。}。複数ページでも脚注が出る。
\end{document}
